\documentclass[12pt]{article}
\usepackage[utf8]{inputenc}
\usepackage[T1]{fontenc}
\usepackage{hyperref}
\title{Dogos\'e Language Resources for K-12 Curriculum Development}
\author{Compiled by Claude Sonnet 5.5 (claude-sonnet-5-5) for Andr\'as Kornai}
\date{}
\begin{document}
\maketitle

\section{Introduction}
Dogos\'e (ISO 639-3 \texttt{dos}; older spelling Doghose) is a Gur language of Burkina Faso with about 20,000 speakers (1991 figure). It is written in Latin script according to one source, and is distinct from Dogoso (\texttt{dgs}). Searches found only general reference pages and no educational materials.

\section{Primary Educational Resources}
No primers, textbooks or school programmes for Dogos\'e were found in search results.

\section{Linguistic and Reference Materials}
\begin{itemize}
\item Overview of the language: \url{https://en.wikipedia.org/wiki/Doghose_language}
\end{itemize}

\section{Religious and Community Resources}
\begin{itemize}
\item Joshua Project profile of the Dogose people: \url{https://joshuaproject.net/people_groups/11588}
\end{itemize}

\section{Digital and Audio Resources}
\begin{itemize}
\item Wikitongues language page: \url{https://wikitongues.org/languages/dos/}
\end{itemize}

\section{Curriculum Development Context}
The language is not known to be taught in schools. Nothing found supports a ready-made curriculum; development would need to start from community-led materials work, taking care not to confuse Dogos\'e with Dogoso.

\end{document}